%% Dans ce fichier, les �ventuels accents sont cod�s
%% en ISO-8859-15 (aka latin9)
%% La ligne suivante donne trois exemples 
%% e accent aigu : �, e accent grave : �, e-dans-l'o : �
%%%%%ENTREES%%%%
\begin{tabarg}
\multicolumn{3}{l}{\titretab {Entr\'ees}}\\
\hline
$X_{in}$ & Vecteur, \texttt{float} & Variable d'origine\\
$toler_{mnq}$ & Coeff., \texttt{Int\_32} & Seuil de calcul du minimum (nb. �chant). A chaque instant de la base de temps d'arriv�e, 
le minimum est calcul� si le nombre d'�chantillons dans l'intervalle de temps est strictement sup�rieur � $toler_{mnq}$. Sinon, 
le r�sultat est "valeur manquante". Si $toler_{mnq}$ = 0, le calcul est fait d�s qu'il y a au moins une valeur valide dans l'�chantillon. \\
$time_{out}$ & Vecteur, \texttt{double} & Base de temps d'arriv�e (secondes)\\
\end{tabarg}

\medskip
%%%%%SORTIES%%%%
\begin{tabarg}
\multicolumn{3}{l}{\titretab {Sorties}}\\
\hline
$X_{out}$ & Vecteur, \texttt{float} & $X_{in}$ minimum �chantillonn� sur
la base de $time_{out}$ (m�me unit� que $X_{in})$\\
\end{tabarg}

